% Expository adaptation, October 7, 2026. Original mathematics: OpenAI, September 23, 2026.
% Prose revised and explanations expanded; see README.md for the source and scope.
\section{Estimating by successive accuracy improvements}
\label{sec:est}

We now supply the averaging schedule described in Section~\ref{sec:estimator-description}. The exact sparse-table identities tell us what to average. The harder question is how accurate each recursive input must be, and how much averaging to apply to it.

The plan is to use estimates at successive child budgets and average their differences. Before giving the schedule, we calculate the cost of changing a bin's scale or moving the source of a query. Those costs determine the precision assignment. We then prove separate bounds for the error in the conditional mean and the noise left by a finite averaging walk.

The prefixes analyzed here are intermediate raw arrays. The full row estimator applies compression after each group and the sparse-table functions after the final group. Its complete evaluation order was given in Section~\ref{sec:estimator-description}; the destination-grouping rule is specified in Section~\ref{sec:fp-blocks}. Here we first analyze the averaging part with exact destination grouping.

The scale and source-domain comparisons first determine the accuracy needed
from each child.  We then average differences between consecutive child
budgets.  Their conditional means telescope, while their small absolute
norms allow finite conditional walks to suppress the noise.  The resulting
prefix estimates are the input to the compression analysis in
Section~\ref{sec:buf-compression}.  That section proves an error reserve
under exact endpoint grouping; Section~\ref{sec:fp-blocks} constructs the
comparison calculation to which the reserve will apply.

\subsection{Tasks and dependency order}

A row task specifies its input matrix, stage, iteration label, start rate
$u$, end scale $a$, optional usage rate $t\le a$, channel roles, and fixed
padding suffixes.  These parameters are common to its whole table.  A
vector task specifies the analogous vertex domain.  Its argument is a
vertex, rather than an additional parameter retained by another task.
Each task has a nonnegative integer statistical budget $M$; $m$ will
denote a child budget.  These budgets are distinct from the numerical
digit places introduced in Section~\ref{sec:arith}.
The ordinary and root discrepancy norms are defined in
Equations~\eqref{eq:atom-row-norm} and~\eqref{eq:atom-vertex-norm}.
We always take $s=8$.

The source weight $X_u(x)/u$ has bounded expectation, and the factor
$1/N$ normalizes the sum over source vertices.  Endpoint sums remain
explicit: bounded row supports or incoming degrees, as appropriate,
will control them.  A bound at one specified start consequently costs
a factor $N$ in moment power; Lemma~\ref{lem:final-majority} makes and
pays this conversion at the end of the proof.

Here is a dependency order, with only a fixed number of task types at each
stage.  First come the bin tasks for $C_l$.  At each pilot step come its
input-score task $\mathcal S_i$, the detector vector $U_i$, and then the
pilot correction task.  Next come the actual correction tasks in iteration
order, the routed bin tasks for $D_l$, and the reward update.  Place the
unnormalized preparation immediately before each normalized task of the
same input.  The raw first channels of both these tasks refer only to
earlier inputs in this order; the normalized task has in addition its
preparation as the source of its second channel.

For a correction task or its preparation, first complete its raw arrays
and pre-incoming scores, then form its incoming gate from those scores.
This assembly uses the current task's completed arrays; every recursive
input to their raw calculation is an earlier task in the order just
specified.

\subsection{Which expressions are averaged}

There are only a few sample formulas to check: a direct entry, a correction followed by a bin, a copy projection, a preparatory mean, a detector trial, and a reward update. The retention-probability ratios below perform changes of conditioning. They are needed even though the nominal rates are comparable to the actual probabilities; an approximate ratio would lose the exact mean identity.

For a fixed row and endpoint let
\[
 \mathcal C_{xy}^{u,a}=\{X_u(x)=1,\ Y_a(y)=1\}.
\]
We use conditional means and conditional walks on this slice.  Write
$\mathcal B_b$ for a normalized bin table and $\mathcal B_b^{(j)}$ for a
bin with destination routing fraction already included; $\mathcal F$
denotes a normalized correction table.  Their values are in atom units
and are zero off their specified membership domains.  In all formulas
below put $\tau=\min(a,b)$.

The direct dense and correction--dense product terms, before division by
the current first-channel scale $a$, are respectively
\begin{align}
 T_b(x,y)&=b\frac{\pi_a}{\pi_\tau}\mathcal B_b(x,y),
 \label{eq:est-direct}\\
 T_b^{\rm prod}(x,y)&=
 b\frac{\pi_a}{\pi_\tau}\frac f{\pi_f}
 \sum_{z:Z_f(z)=1}\mathcal F(x,z)\mathcal B_b(z,y).
 \label{eq:est-product}
\end{align}
In the product, the left task has roles $X,Z$ and rates
$(u;f,f)$, listed as start, end, and usage.  The right task has roles
$Z,Y$ and rates $(f;b,\tau)$.  In the direct term the start rate stays $u$.
The term $-E$ in $D$ is the negative of
Equation~\eqref{eq:est-direct} with $b=f$ and $\mathcal F$ in place of
$\mathcal B_b$.  Summing the direct and product terms gives the raw inputs
$G_i,G_i'$; adding the negative term gives $D$.

For a copy projection $x_i,y_j$ from stage $l+1$ to stage $l$, the raw term
is
\begin{equation}
 r_i(x)b\frac{\pi_a}{\pi_\tau}\mathcal B_b^{(j)}(x,y).
 \label{eq:est-projection}
\end{equation}
Partition by the fixed symbols $i,j$.  Append those symbols to the native
hash arguments before the inherited suffixes on their respective channels.
The parent-stage source multiplier is obtained from its detector vectors
on the broad parent-stage vertex domain.  No child tuple contains the old
source or a particular right endpoint as an extra parameter.

For the second raw channel, use $\zeta$ times the preparatory table, take
its conditional average on $\mathcal C_{xy}^{u,t}$, and extend the result
by zero off $Y_t(y)$.  For a detector vector use the trial
\begin{equation}
 \operatorname{rank}_{J}
       \{\mathcal S_i(x,y):X_f(x)=1\}
 \quad\hbox{on the slice }Y_f(y)=1.
 \label{eq:est-detector}
\end{equation}
Ranks include arbitrary zero padding.  For rewards, writing $f=b_l$, the
two terms for $W_{l+1}(x_i)$ on its inherited source slice are
\begin{equation}
 r_i(x)W_l(x),\qquad
 r_i(x)\frac f{\pi_f}
       \sum_{z:Z_f(z)=1}\mathcal F(x,z)W_l(z).
 \label{eq:est-reward}
\end{equation}
The first $W_l$ call keeps the inherited source domain; the second restarts
on the midpoint domain.  Vector answers are finally clipped into $[0,1]$.

\begin{lemma}[True sample means]\label{lem:est-true-means}
Using true stochastic tables in
Equations~\eqref{eq:est-direct}--\eqref{eq:est-reward} gives, after the
indicated conditional averages, the prescribed raw inputs and vector
values on their conditioning slices, with the stated zero extensions.
In particular, the first row channel has target
$d_y=G(x,y)/a$ and the second has target $\zeta q_y\mathbf1_{Y_t(y)}$.
\end{lemma}
\begin{proof}
First consider how the destination's retention condition changes with scale. If $b\ge a$, the usage hit is already the conditioning hit.  The
conditional normalization of a bin therefore turns $b\mathcal B_b$ into
its deterministic bin contribution.  If $b<a$, nested origin boxes give
\[
 \mathbb E\left[\frac{\pi_a}{\pi_b}\mathbf1_{Y_b(y)}V
                   \,\middle|\,\mathcal C_{xy}^{u,a}\right]
   =\mathbb E[V\mid\mathcal C_{xy}^{u,b}]
\]
for every integrable $V$ on the narrower slice.  Thus the importance factor
performs exactly the required change of conditioning.

For a product, we must justify separately why the two true normalizations can be used. For a fixed midpoint in Equation~\eqref{eq:est-product}, first make this
terminal change when necessary and then condition on $Z_f(z)=1$.
Apart from membership indicators, the true right bin gate and denominator
use only channel $Y$.  The true left correction uses only channels $X,Z$.
These channels are independent, and the remaining true multipliers are
deterministic.  The two normalizations can therefore be applied separately.
In particular, after the midpoint hit is fixed, the shared role $Z$
does not occur in the right bin's remaining random value.  This is an
identity for the true atoms; the estimates will instead be compared by
their joint error norms.
The factor $f/\pi_f$ cancels midpoint sampling and converts its correction
atom to the correction entry.  Summing midpoints proves the product
formula.  The projection multiplies by deterministic $r_i(x)$ and includes
the routing fraction in the bin, so the same argument applies on each
fixed-symbol part.  Fixed-suffix invariance in Lemma~\ref{lem:env-ranks}
justifies all inherited padding conventions.

The second channel is its defining preparatory conditional mean, on its
own narrower slice; it has no importance factor.  The detector mean is
the definition of $U_i$.  True reward factors are deterministic, so the
same midpoint argument proves Equation~\eqref{eq:est-reward}.
\end{proof}

To form a budget-$m$ sample term, replace all earlier inputs in the
corresponding formula by budget-$m$ estimates, with their domains and
common task parameters unchanged.  This includes the detector inputs
to $r_i$; the primary factor $r_0=1$ remains constant.  Products, source
functions, and ranks are evaluated before averaging.  We first compare
this sample's conditional mean with the same formula using true atoms.
The scale and source-domain costs of that comparison will determine
how much accuracy a budget must provide.

\subsection{Mean comparisons and the cost of releasing a source}

Lemma~\ref{lem:est-true-means} identifies the exact conditional means.
For estimated factors, the comparison has two distinct costs.  The
factor $b/a$ converts a bin contribution to the current entry units.
A right factor at a midpoint also replaces the old source domain of
rate $u$ by the stage domain of rate $f$.  The following estimates
account for both changes without assuming independence of the errors.

The following scalar loss packages the comparison in a form that behaves well under averaging. Its denominator discounts an error in proportion to the magnitude of the values being compared. It is equivalent to the eighth-root discrepancy already used for the sparse tables, but its homogeneity makes the scale calculation more transparent.

For nonnegative $H,H'$ define
\begin{equation}
 \mathcal L(H,H')=
 \begin{cases}|H-H'|^s/(H+H')^{s-1},&H+H'>0,\\0,&H=H'=0.
 \end{cases}
 \label{eq:est-loss}
\end{equation}

\begin{lemma}[Homogeneous loss]\label{lem:est-loss}
The loss in Equation~\eqref{eq:est-loss} is comparable, by constants
depending only on $s$, to $|H^{1/s}-(H')^{1/s}|^s$.  It contracts under
sums and nonnegative mixtures.  If $H\le C_0G$ for a nonnegative $G$, then
for every $a>0$ and $H'\ge0$,
\begin{equation}
 a\min\left\{1,\frac{|H-H'|}{a+G}\right\}^{s}
       \le C\mathcal L(H,H'),
 \label{eq:est-relative-loss}
\end{equation}
where $C$ depends only on $s,C_0$.
\end{lemma}
\begin{proof}
Put $v=\max(H^{1/s},(H')^{1/s})$ and
$w=\min(H^{1/s},(H')^{1/s})$.  Factoring $v^s-w^s$ bounds it between
$v^{s-1}(v-w)$ and $sv^{s-1}(v-w)$; the denominator is between
$v^{s(s-1)}$ and $2^{s-1}v^{s(s-1)}$.  This proves comparability.
For a mixture, apply H\"older to
$|H-H'|=\mathcal L(H,H')^{1/s}(H+H')^{(s-1)/s}$ and then integrate;
the triangle inequality first bounds the absolute difference of the
integrals.  This proves contraction, including finite sums.

If $H'\le2H+a$, then $H+H'\le3C_0G+a\le C'(a+G)$, which proves
Equation~\eqref{eq:est-relative-loss} directly.  If $H'>2H+a$, then
$|H-H'|>H'/2$, $H+H'<3H'/2$, and $H'>a$.
Thus $\mathcal L(H,H')\ge cH'\ge ca$, while the left side is at most $a$.
\end{proof}

Every unsigned true term of the signed formula for $D$ is bounded by a
constant times $D$, by Lemma~\ref{lem:hier-residual}; the corresponding
assertion for the other raw inputs follows from their nonnegative sums.
Consequently the preceding lemma can be applied separately to all true
contributions.  Triangle inequalities combine their clipped relative
errors even when a contribution has a negative sign.

\begin{lemma}[Scale and domain comparison]\label{lem:est-mean-comparison}
Compare a budget-$m$ sample term with the same term using true stochastic
inputs.  For a first-channel mean contribution use its clipped relative
error against $1+d_y$.  In the bound for its $s$th-power norm, the scale
factor multiplying a child root-discrepancy power is, up to fixed constants,
\[
 \begin{cases}
 b/a,&b\ge a,\\
 (b/a)^{s-1},&b<a.
 \end{cases}
\]
A child that releases the old source and restarts on a stage-rate $f$
domain costs in addition at most $f/u$ in norm power.  Direct children
retaining that source have no such factor.  The same-scale vector and
second-channel comparisons have constant scale cost, with $f/u$ only
where a source is released.  No independence of estimated factors is
required.
\end{lemma}
\begin{proof}
We treat coarse and fine bins separately because changing the conditioning contributes a different power of the scale ratio in the two cases. We then check vectors and preparatory means.

First suppose $b\ge a$.  Apply
Equation~\eqref{eq:est-relative-loss} after conditional expectation, and
then the mixture contraction of Lemma~\ref{lem:est-loss} before it.
Homogeneity supplies a factor $b/a$ after division by $a$.
The factor $f/\pi_f$ in a product is bounded above and below by fixed
constants.  For bounded nonnegative atoms a product-root difference splits
into the factor-root differences, multiplied by fixed constants times
support indicators of the unaffected factor in one of the two versions.

The source-domain factor is a counting issue, not an independence claim. An error at a midpoint can be reached from only boundedly many old sources through a correction access. To determine this factor, write
$\Delta_s\mathcal F$ and $\Delta_s\mathcal B_b$ for the root discrepancies.
After integration on the source/end slice, the relevant sums, besides
$O(b/a)/N$, have the forms
\begin{align*}
 &\mathbb E\sum_{x,z,y}\frac{X_u(x)}u Z_f(z)Y_a(y)
       |\Delta_s\mathcal F(x,z)|^s
       \mathbf1_{\mathcal B_b(z,y)\ne0},\\
 &\mathbb E\sum_{x,z,y}\frac{X_u(x)}u Z_f(z)Y_a(y)
       \mathbf1_{\mathcal F(x,z)\ne0}
       |\Delta_s\mathcal B_b(z,y)|^s.
\end{align*}
Either support indicator can come from either compared version.  In the
first sum the right row cap eliminates $y$ at constant cost.  In the
second sum the left incoming cap eliminates $x$ at cost $O(1/u)$.
The right child's row metric has weight $1/f$, leaving exactly $O(f/u)$.
Its tuple is the same for every old source, so this elimination compares
the same right error throughout.  More explicitly, for either left
support relation $\mathcal P$ with incoming degree at most $K_c$ and
a common right discrepancy $e(z,y)$, the relevant inequality is
\begin{align*}
 &\frac1N\mathbb E\sum_x\frac{X_u(x)}u
       \sum_{z:(x,z)\in\mathcal P}Z_f(z)
       \sum_{y:Y_b(y)=1}|e(z,y)|^s\\
 &\qquad\leq \frac{K_c}{Nu}\mathbb E\sum_z Z_f(z)
       \sum_{y:Y_b(y)=1}|e(z,y)|^s\\
 &\qquad=K_c\frac fu\,\|e\|_{f,b;s}^s,
\end{align*}
where the source channel in the last norm is $Z$.
The incoming cap applies pointwise, so the discrepancy can depend on
the same environment as $\mathcal P$.
Using $Y_a\le Y_b$ enlarges the remaining
endpoint domain to the child's own domain.  A direct term simply retains
its source and endpoint sums.  In a projected term, the source-multiplier
error eliminates $y$ by a row cap and changes the vertex weight from
$1/u$ to $1/f_{\rm par}$; this costs $O(f_{\rm par}/u)$.
Partitioning the fixed copy symbols keeps every tuple common during these
sums.

For fine bins, using the coarse-bin estimate would miss the useful small factor. For $b<a$, change the terminal conditioning to $Y_b$ \emph{before}
taking a power.  The normalized coefficient inside is then $b/a$.
There are only constantly many nonzero left paths in the union of the
two rows.  Ordinary bounded-factor difference inequalities and the
finite-sum power inequality therefore apply with constant cost.  The
integrated change from $Y_a$ to $Y_b$ leaves $\pi_a/\pi_b$ outside the
power.  Its total factor is
\[
 (b/a)^s\frac{\pi_a}{\pi_b}=O((b/a)^{s-1}).
\]
The same row and incoming sums as above then apply, now with $Y_b$.
Ordinary differences of bounded nonnegative values are controlled by
their root differences.  Absolute mean error bounds clipped relative
mean error, so this proves the fine-bin assertion.

For the second raw channel, Jensen is applied directly on
$\mathcal C_{xy}^{u,t}$.  Its integrated endpoint mask is $Y_t$, already
contained in the preparatory norm's $Y_a$ domain.  There is no importance
factor and hence no $\pi_a/\pi_t$ penalty.  The fixed scaling $\zeta$
does not increase the cost.

A rank discrepancy is at most the supremum of its input discrepancies.
For the detector trial, summing its power over $Y_f/f$ is bounded by the
score error sum over $X_f/f$ and $Y_f$; the rates match exactly.
For rewards, the finite left-path count gives an ordinary product
comparison, and the same incoming cap eliminates old starts for a
right-vector error.  Source functions have bounded Lipschitz roots in
their clipped detector inputs, so their costs are the broad vertex costs
already computed.  Every estimate here is pointwise or follows from
Jensen and summation; none uses independence of its errors.
\end{proof}

The preceding comparison measures errors in a retained sample term.
Omitting a fine bin instead discards its whole true contribution; the
next bound measures that loss.

\begin{lemma}[Fine omission]\label{lem:est-omission}
If one fine bin is omitted, the norm of its normalized contribution is at
most $C(b/a)^{(s-1)/s}$.  This applies to direct terms, routed terms, and
correction--bin products.
\end{lemma}
\begin{proof}
The deterministic contribution $H$ has entries at most $Cb$ and bounded
row mass.  For a product these bounds follow from
Lemma~\ref{lem:hier-bin-products}.  Therefore
\[
 \sum_y\pi_a|H(x,y)/a|^s
 \le \frac{\pi_a}{a^s}(Cb)^{s-1}\sum_yH(x,y)
 \le C'(b/a)^{s-1}.
\]
Integrating source hits gives $\pi_u/u=O(1)$, and summing at most $N$ rows
cancels the normalization $1/N$.
\end{proof}

\subsection{Choosing precision to pay for scale and domain changes}

The preceding bounds explain the rate weights in the accuracy parameter.
Write $h_b=\log_2(b/a)$ for the offset of a bin scale.
In norm, the coarse-bin cost is $(b/a)^{1/s}$ and the fine-bin cost is
$(b/a)^{(s-1)/s}$.  Since $s=8$, a weight $1/2$ between these two
exponents leaves a geometric margin in both directions.  The same
weight pays for the source-release cost $(f/u)^{1/s}$.  We now make
this accuracy convention precise; the additive schedule will then
trade lower child budgets for longer averaging walks.

Assign positive integer ranks $i_*$ in the task order fixed above, so every
dependency decreases rank by at least one.  The finite number of pilot steps and of
preparations is an absolute constant, so $i_*=O(l+1)$.  A task at integer
budget $M\ge0$ has implicit precision
\begin{equation}
 p=\frac{M}{K_{\rm prec}}-\frac12\log_2\frac1a
        -\frac12\log_2\frac1u-A_*i_*.
 \label{eq:est-precision}
\end{equation}
For a stage-$l$ vector use $a=f=b_l$; detector vectors also have $u=f$,
whereas rewards retain their inherited start rate.  The constants
$K_{\rm prec}$ and $A_*$ will be chosen below.  If $p\le0$, return the zero
table or vector and empty access classes.  Otherwise the task is called
active.  The intended induction bound is $C_b2^{-p}$ for every required
ordinary vector norm and every required root discrepancy norm, including
the pre-incoming scores and witnesses.  A fixed $C_b$ makes this assertion
trivial for inactive tasks: values and row supports are uniformly bounded,
and $\mathbb E X_u/u=\pi_u/u=O(1)$.

At a normalized correction task, the column multiplier is assembled from
earlier detector-vector estimates at budget $M-c_{\rm col}$, for a fixed
positive integer $c_{\rm col}$.  This is a separate final assembly call.
All other calls used in one additive sample term have the common
budget denoted by $m$.

Suppose a child at budget $m$ has parameters $a',u',i_*'$.  Subtracting
Equation~\eqref{eq:est-precision} for the parent gives the exact identity
\begin{equation}
 p'-p=-\frac{M-m}{K_{\rm prec}}
       +\frac12\log_2\frac{a'}a
       +\frac12\log_2\frac{u'}u
       +A_*(i_*-i_*').
 \label{eq:est-child-precision}
\end{equation}
In a bin comparison $a'\ge c b$ for a fixed $c>0$: a bin child has
$a'=b$, while a correction has $a'=f$ and $b\le K_f f$.
The start is either retained or reset to the indicated broad stage rate.
Fixed-stage changes in vector calls only change fixed constants.
Thus, after taking $s$-th roots of the costs in
Lemma~\ref{lem:est-mean-comparison}, the scale contribution times child
accuracy is bounded by
\begin{equation}
 C C_b\,2^{-p-A_*+(M-m)/K_{\rm prec}}
          2^{-3|h_b|/8}.
 \label{eq:est-scale-slack}
\end{equation}
Indeed for $h_b\ge0$ the exponent is $h_b/s-h_b/2=-3h_b/8$;
for $h_b<0$ it is $(s-1)h_b/s-h_b/2=3h_b/8$.
For a child that restarts at rate $f$, the source-domain comparison costs
$(f/u)^{1/s}$ in norm, while its precision gains
$\frac12\log_2(f/u)$.  Multiplying the cost by the resulting accuracy
factor gives
$(f/u)^{1/s-1/2}=(f/u)^{-3/8}\le C$, using the allowed start-rate
condition $u\le f$.

For the separate final column-vector call, $a=f$, its child uses
$a'=u'=f$, and $M-m=c_{\rm col}$.  Thus its norm cost times child
accuracy is at most
\[
 C C_b2^{-p}\,2^{c_{\rm col}/K_{\rm prec}-A_*}
       (f/u)^{-3/8}.
\]
The fixed rank separation pays this fixed budget loss.  This calculation
also explains why the source-domain ratio is charged in the error
analysis rather than inserted as a numerical coefficient in a sample.

\subsection{Average differences and group them by budget decrease}

We can now choose the schedule using the precision comparison just proved. A direct high-budget child call would save little budget for the caller's own work. Instead, express a high-budget sample as a sum of successive improvements. Low-budget improvements can tolerate long averaging walks; later improvements are smaller because their two child estimates are more accurate.

Every budget-zero task is inactive under the preceding precision rule.
Hence a budget-zero sample term is zero: it contains at least one
estimated table or vector factor.  Exact base terms $S(x,y)/a$ and
$W_0(x)$ are inserted separately.

For every first-channel term labelled by scale $b$, set
$c_b=\lceil c_0(1+|h_b|)\rceil$, where $c_0$ is a fixed positive constant.
This rule includes the sparse direct term $-E$ with $b=f$, as well as
direct bins, products, and projection terms.  Every child factor in a
term, including the detector-vector inputs to a projected source
multiplier, uses that term's common budget $m$ and cutoff $c_b$.
Use zero scale offset only for second-channel preparatory means,
detector-vector trials, and reward terms with fixed stage-rate changes.
The separate final column assembly retains its fixed drop
$c_{\rm col}$.  Denote the corresponding normalized
sample term at budget $m$ by $T_{b,m}$, including division by $a$ for the
first row channel.  For $1\le m\le M-c_b$ form
\begin{equation}
 T_{b,m}-T_{b,m-1}.
 \label{eq:est-difference}
\end{equation}
Apply to this difference the conditional walk of length
$\ell(M-m)$, where $\ell$ is a sufficiently large fixed integer.  Use the
same generator word and the same guessed wraps for its two parts.

We call $M-m$ the slack of a child call: it is the amount by which its statistical budget has decreased. The same decrease will later pay for the description of the call. Group increments by dyadic slack
$d\le M-m<2d$.  Add groups in decreasing powers of two $d\le M$ down to
$1$, including empty groups.  Put exact base contributions in the first
group.  Let $S_j$ be the uncompressed prefix through the group with slack
$d_j$.  The conditional mean prefix is obtained by replacing each walk by
uniform averaging on its slice; noise is the difference from that mean.
For an available bin the last budget reached by this mean prefix is
\begin{equation}
 m_j(b)=\max\{0,M-\max(d_j,c_b)\}.
 \label{eq:est-last-budget}
\end{equation}
This exact cutoff will matter in the error bound.  Bins with $c_b\ge M$
have no layers and are omitted.


The walk is applied to an array indexed by actual endpoints.  Thus,
for each term and pair, its input is the single function
\[
 F_{xy}(\sigma)=T_{b,m}(x,y;\sigma)-T_{b,m-1}(x,y;\sigma),
\]
with all contributions at that endpoint already summed.  A path
description is only a way to enumerate those contributions.  In
particular, the two budgets must be grouped at the same actual
endpoint before their difference is measured.  The following lemma
establishes this identity for the path enumeration and also records
the description length needed later by the controller.

\begin{lemma}[Immediate paths]\label{lem:est-paths}
At a row group of slack $d$, immediate candidates have descriptions of
$O(d)$ bits.  With exact endpoint grouping their sums equal the specified
conditional-walk increments.  A product visits its midpoint by a paired
correction move before using any factor at that midpoint.  Its right-hand
task tuple does not depend on the discarded old source.  Every child
estimator in such a candidate has budget at most $M-d$.
\end{lemma}
\begin{proof}
We must check both what a path describes and why summing those paths gives the intended average. A descriptor records the finitely many formula and copy-part choices, its
scale and budget offsets, its sign in Equation~\eqref{eq:est-difference},
the generator word and guessed wraps, and a bounded number of constituent
ports.  An included scale obeys $|h_b|=O(d)$, and the walk has length
$O(d)$.  Lemma~\ref{lem:env-variants} therefore bounds all these records by
$O(d)$ bits.  Port alphabets are fixed by Lemma~\ref{lem:atom-ports}.
There are only $O((1+d)^2)$ choices of budgets and scales before the
exponential word/variant choices.

Check each constituent port, the outer source and endpoint memberships,
and the guessed wraps at the vertices where their coordinates can be
evaluated.  For a fixed endpoint and original generator word, unique wrap
validation leaves exactly its genuine conditional transformation.
Give the surviving term its original walk-word probability, multiplied
by the formula coefficients and sign.  Guessed wraps do not change that
probability.  Constituent ports enumerate entries once.  Hence summing
over true equal endpoints is precisely the walk of the difference, even
though its endpoint was not known when the variant was specified.

The first factor of every product is a correction, whose retained ports
have bounded degree in both directions.  Its paired move supplies the
midpoint and a bounded return token.  The right factor can then be queried
using the common $Z,Y$ tuple prescribed above; that tuple contains no old
source.  A direct unbalanced terminal motion has no further factor query
at its endpoint.  These are combinatorial access requirements; the
controller section implements the pairing and terminal operations.
Finally, the slack definition gives $m\le M-d$, also for the $m-1$ term.
\end{proof}


All reachable rate logarithms remain $O(B)$ when $M_{\rm top}=O(B)$:
a scale extension at slack $d$ changes a reciprocal-rate logarithm by
$O(d)$ and reduces budget by at least $d$.  Stage resets have only their
fixed stage rates, usage rates are minima of available rates, and inherited
padding has at most one symbol per stage above the task.  Summing along a
dependency chain bounds all extensions by $O(M_{\rm top})$.  The field in
the shared environment can therefore be chosen large enough for every
used scale, including budget-zero endpoints of available differences.
Omitted finer deterministic contributions need no hash tables.


We can now estimate each mean prefix by telescoping the conditional
means to budget $m_j(b)$.  Its deterministic omission error and its
child-estimation error are treated separately.  The remaining noise
will be bounded by applying the conditional spectral estimate to the
endpoint functions $F_{xy}$ above.

\begin{lemma}[Mean prefixes]\label{lem:est-mean-prefix}
Assume all earlier table and vector bounds with coefficient $C_b$.
For an active task there is a fixed $B_1$, independent of the later
choice of $A_*$, such that every mean prefix has norm at most
$c_*2^{-p+B_1d_j}$ against its target.  For the first channel the norm
uses clipped relative discrepancy; for the second and for vectors it uses
ordinary absolute discrepancy.  The coefficient $c_*$ can be made
arbitrarily small by increasing $A_*$.
\end{lemma}
\begin{proof}
There are two sources of mean error. At a scale included in the schedule, telescoping stops at a finite child budget. At an omitted scale, the entire deterministic contribution is missing. We bound them separately.

\emph{Available scales.}  For an available scale, conditional expectations of
Equation~\eqref{eq:est-difference} telescope to its term at
$m=m_j(b)$ in Equation~\eqref{eq:est-last-budget}.  If this budget is
zero, the term is zero and the inactive-child bound still applies.
More generally, the child bound remains valid whenever its precision is
nonpositive, by the bounded-output estimate for inactive tasks.
In every case
\[
 M-m\le d_j+c_0(1+|h_b|)+1.
\]
Choose $K_{\rm prec}$ so that $c_0/K_{\rm prec}<1/8$.
Equation~\eqref{eq:est-scale-slack} is then bounded by
\[
 C' C_b\,2^{-p-A_*+d_j/K_{\rm prec}}2^{-|h_b|/4}.
\]
The rate grid is geometric, so summing this bound over available scales
costs a fixed constant.  There are only a fixed number of formula and
factor types.  Unbinned calls use the same calculation without the scale
sum.  Minkowski and subadditivity of clipped relative discrepancies
combine all these terms.  Exact base contributions are already present
in every prefix and have no error.

\emph{Wholly omitted scales.}  A bin with $c_b\ge M$ contributes no
layer.  We bound its deterministic contribution directly by
Lemma~\ref{lem:est-omission}.  First we show that every such bin is fine.
Put $r=\log_2(1/a)\ge0$.
A coarse bin has $0\le h_b\le r$ because all grid rates are at most one.
Activity gives
$M>K_{\rm prec}(r/2+A_*i_*)$, even after discarding the nonnegative
start-rate cost.  With $K_{\rm prec}/2>c_0$ and $A_*$ sufficiently large,
this implies $M>\lceil c_0(1+r)\rceil$, so a wholly omitted bin is fine.
For the sparse direct term $b=f$, the same argument gives availability
when $f\ge a$.  When $f<a$, the admissibility condition $a\le K_f f$
gives $|h_f|\le\log_2 K_f$, and increasing $A_*$ ensures $M>c_f$.
Thus this term is never wholly omitted in an active task.
The inequality $c_b\ge M$ implies
\[
 |h_b|\ge (M-1)/c_0-1.
\]
Lemma~\ref{lem:est-omission} and a geometric sum consequently bound all
such omissions by
$C2^{-(s-1)M/(s c_0)}$.  Since $p\le M/K_{\rm prec}$ and
$K_{\rm prec}>s c_0/(s-1)$, this is at most
$C2^{-p-\eta M}$ for some fixed $\eta>0$, after absorbing the fixed
rounding offset.  Finally $M>K_{\rm prec}A_*i_*\ge K_{\rm prec}A_*$
makes its coefficient arbitrarily small as $A_*$ increases.
Taking, for example, any fixed $B_1>1/K_{\rm prec}$ proves the claim,
including every last-budget and no-budget cutoff loss.
\end{proof}

\subsection{Absolute difference bounds and averaging noise}

The conditional-mean errors are now bounded.  To control the noise left
by finite averaging, we need an absolute norm bound on each paired
difference before the walk is applied.  This coarser bound will also be
used for the fingerprint comparison.

\begin{lemma}[Absolute differences]\label{lem:est-raw-difference}
At a slack-$d$ group, the ordinary raw norm of an unaveraged paired
difference is at most
\begin{equation}
 C C_b\,2^{-p-A_*+C_0d},
 \label{eq:est-raw-bound}
\end{equation}
where $C_0$ is fixed independently of the walk length coefficient.
For two coupled versions of a budget-$r$ sample term, if the
corresponding child discrepancies are bounded by
$2^{-p'_r}\lambda_r$, their normalized raw difference is bounded by
$2^{-p+C_0'd}\lambda_r$, summed over the finitely many child types.
For a paired layer at budgets $m$ and $m-1$, apply this estimate to
both parts. Its coupled discrepancy is consequently bounded by
$2^{-p+C_0'd}(\lambda_m+\lambda_{m-1})$, after enlarging the fixed
constant, or by the corresponding maximum over all children in the
group. Every child in either part has budget at most $M-d$.
For a fixed guessed measure-preserving variant and its valid endpoint
restriction the same bounds hold.  Counting all variants separately may
increase $C_0'$ by a constant depending on the chosen walk length.
\end{lemma}
\begin{proof}
The useful cancellation comes from comparing the two child budgets with one common true sample, not by bounding their magnitudes separately. Subtract the \emph{same stochastic true-atom term} from the budget-$m$
and budget-$(m-1)$ terms, and use the triangle inequality.  A product has
only constantly many left paths in the union of the compared rows.
Ordinary bounded-factor differences and the row and incoming caps
therefore give the sums used in
Lemma~\ref{lem:est-mean-comparison}, without any averaging.  For a coarse
bin the absolute coefficient costs $b/a$ in the norm rather than its
$s$-th root.  For a fine bin the importance factor can stay inside the
power.  Both are bounded by $2^{O(|h_b|)}$.  Equation~\eqref{eq:est-child-precision}
pays the only separate possibly large domain ratio, $f/u$, as before.
All other scale and accuracy changes are $2^{O(d)}$ because
$|h_b|=O(d)$ and $M-m=O(d)$ in this group.  Ranks and bounded source
functions obey the same pointwise comparisons.  This proves
Equation~\eqref{eq:est-raw-bound}.  Replacing comparison to truth by a
comparison of the two coupled child versions gives the second assertion
by the identical inequalities; independence is still unnecessary.

For a fixed guessed variant, valid paths require the transformed
memberships.  Change variables by its full-environment measure-preserving
bijection and then discard surplus validity guards.  This leaves exactly
the child-domain sums just bounded.  When combining the paired terms,
combine their contributions at the same actual endpoint before measuring
the difference.  The number of distinct variant descriptors is
$2^{O(d)}$, so treating them separately merely enlarges the crude exponent.
The estimate before such separate counting is a statement about the
single-seed sample function and has no dependence on the eventual walk
length coefficient.
\end{proof}

\begin{lemma}[Conditional prefix bounds]\label{lem:est-prefix}
Assume exact local grouping, the earlier task error bounds, and that the
sample functions depend only on variables varied by the conditional walk.
Their uniform slice means are then deterministic at each indicated pair
or vertex.  Write $S_j^{\rm mean}$ for a mean prefix and
$S_j^{\rm noise}=S_j-S_j^{\rm mean}$ for its noise.  There are fixed
$B_1$ and arbitrarily large prescribed $K_n$ such that
\begin{align}
 \left\|\min\left(1,
       \frac{|S_j^{1,\rm mean}-d|}{1+d}\right)\right\|_{u,a;s}
       &\le c_*2^{-p+B_1d_j},
       \label{eq:est-prefix-mean}\\
 \|S_j^{\rm noise}\|_{u,a;s}
       &\le c_*2^{-p-K_nd_j}.
       \label{eq:est-prefix-noise}
\end{align}
The second-channel mean uses its ordinary absolute norm, extended by zero
off its usage domain, and satisfies the first bound.  Vector means and
noises satisfy the corresponding vertex bounds.  The coefficient $c_*$
can be chosen as small as needed by task separation and walk constants.
\end{lemma}
\begin{proof}
The mean assertion is Lemma~\ref{lem:est-mean-prefix}. For noise, apply the conditional mixing estimate to each already-grouped destination function, then sum the resulting moments. This order avoids paying separately for every guessed walk variant before obtaining the mixing gain.
For a difference layer, write $F_{xy}$ for its sample function on
$\mathcal C_{xy}^{u,a}$, $\mathcal A_{xy}$ for the conditional walk
operator of length $\ell(M-m)$, and $\Pi_{xy}$ for uniform averaging
on that slice.  Put
$w_{xy}=\Pr(\mathcal C_{xy}^{u,a})/(Nu)$ and
$r_{M-m}=C2^{-c\ell(M-m)}$.
Lemma~\ref{lem:env-gap} gives a uniform operator bound, so
\[
 \sum_{x,y}w_{xy}
   \|(\mathcal A_{xy}-\Pi_{xy})F_{xy}\|_{L_s}^s
 \le r_{M-m}^s
   \sum_{x,y}w_{xy}\|F_{xy}\|_{L_s}^s.
\]
The slice weights make the sum on the right precisely the integrated raw
norm to the $s$th power, bounded by Lemma~\ref{lem:est-raw-difference}.
Thus the uniform operator estimate applies directly to the integrated
input bound, even when different slices have different input norms.
For the second channel replace $Y_a$ by $Y_t$ in the slice and its
weight; the mask stays $Y_t$ throughout.  Vectors use the same
calculation on their one-vertex slices.  No extra conditioning factor
appears.

There are $O((1+d)^2)$ budget/scale terms in a group.  Choose $\ell$ so
that the contraction dominates $2^{C_0d}$ and this polynomial, leaving
at least $2^{-(K_n+1)d}$.  Summing groups with slack at least $d_j$ gives
Equation~\eqref{eq:est-prefix-noise}.  Exact base terms have no noise.
This contraction applies to the paired difference as one function of the
same seed.  The sum of valid variants is the conditional walk, so the
gap applies to that sum before any separate variant counting.
Increasing $A_*$ supplies an
arbitrarily small coefficient in both this estimate and the mean bound.
\end{proof}

The determinism hypothesis in Lemma~\ref{lem:est-prefix} applies at
each fixed pair or vertex, in the task's effective environment as in
Remark~\ref{rem:env-effective}.  Uniform averaging must integrate every
coordinate on which its sample function depends.  In
Section~\ref{sec:fp-blocks}, the clean calculation removes dependence
on the current block's salt, while the parent walk varies all earlier
salts used by lower-block calculations; Lemma~\ref{lem:fp-clean}
verifies this hypothesis.  The resulting first-channel mean is a
deterministic number, and the second-channel mean has only its prescribed
$Y_t$ membership mask.  These are the means used in the conditional
heavy-count argument of Lemma~\ref{lem:buf-rarity}.

